\section{Introdução}
\label{sec:intro}

A classificação atribui categorias a uma imagem. O \emph{retrieval} ordena uma coleção pela relevância de cada imagem para uma categoria ou query. Ambas dependem da representação que transforma a imagem num vetor de dimensão fixa. Comparamos representações de três gerações num protocolo comum e medimos o AP/mAP, a eficiência em labels, a estabilidade entre seeds e o custo.

\paragraph{Representações handcrafted.}
O BoVW \cite{sivic2003,csurka2004} atribui cada descritor SIFT denso \cite{lowe2004} à palavra mais próxima de um vocabulário k-means e representa a imagem pelo histograma das palavras (dimensão $K$). Os tiles espaciais \cite{lazebnik2006} recuperam parte do layout. Sobre estes vetores treina-se
um SVM linear \cite{cortes1995}. A raiz quadrada do histograma em L1 é um mapa explícito exato do kernel de Hellinger \cite{vedaldi2012}. A VLAD \cite{jegou2010} soma, por palavra, os resíduos ao centro (dimensão $KD$). O Fisher Vector \cite{perronnin2010} usa um GMM com atribuição suave e codifica
estatísticas até à 2.ª ordem ($2KD{+}K$). Como estes vetores crescem depressa, perguntamos também se um vocabulário grande comprimido por PCA vale mais do que um vocabulário pequeno com a mesma dimensão.

\paragraph{Representações aprendidas.}
Com o backbone congelado e o mesmo SVM linear, comparamos pré-treino supervisionado (ResNet-18 \cite{he2016resnet}, ViT-B/32
\cite{dosovitskiy2021vit}), auto-supervisionado (DINOv2 \cite{oquab2024dinov2}) e imagem--texto (CLIP \cite{radford2021clip}). O CLIP permite também classificar e pesquisar por texto sem labels (\emph{zero-shot}).

\paragraph{Principais resultados.}
\begin{itemize}\setlength\itemsep{0pt}
    \item \textbf{BoVW (Sec.~\ref{sec:part1}):} Hellinger com tiles $2\times2$ é a melhor combinação. Com 10\,\% do treino, iguala o L2 com 50\,\%.
    \item \textbf{VLAD/FV (Sec.~\ref{sec:part2}):} com $K{=}32$, o mAP de test sobe de 0.694 (BoVW) para 0.746 e 0.784.
    \item \textbf{PCA (Sec.~\ref{sec:part3}):} à mesma dimensão, vocabulário grande + PCA supera vocabulário pequeno, mas o custo de extração não baixa.
    \item \textbf{Pré-treinados (Sec.~\ref{sec:part4}):} ResNet-18 0.964, ViT-B/32 0.961, DINOv2 0.981 (CLS), CLIP + SVM 0.991. Em validação, o DINOv2 com 1 positivo supera em média o BoVW com todos (0.779 contra 0.717). O CLIP zero-shot chega a 0.826.
    \item \textbf{Retrieval por query (Sec.~\ref{sec:part5}):} sem SVM, BoVW e VLAD ficam perto do acaso (mAP 0.105 e 0.119, aleatório 0.060). O DINOv2 chega a 0.806 e o CLIP texto a 0.803.
\end{itemize}
A Sec.~\ref{sec:discussion} junta os resultados das cinco partes e discute custos e limitações. A Sec.~\ref{sec:conclusion} conclui.
